% GENERATED FILE. Edit canonical lessons, then run npm run generate:books.
% canonical-source-hash: fnv1a64:ecfab57d65c414b1
% canonical-lessons: RU-C303-rozhdeniye, RU-C303-rozhdatsya, RU-C303-mladenets, RU-C303-beremennaya, RU-C303-rozhat

\chapter{Birth}
\label{ch:ru-birth}
\hlchaptermodality{303}

\begin{chapteropening}
\textbf{By the end of this chapter:} \emph{I can talk about a birth, being born, a baby, being pregnant and giving birth.}

Complete the last lesson of chapter 303: I can talk about a birth, being born, a baby, being pregnant and giving birth.
\end{chapteropening}

\section[\texorpdfstring{\ru{рождение} (rozhdéniye)}{рождение (rozhdéniye)}]{\ru{рождение} (rozhdéniye) — birth}

\label{lesson:RU-C303-rozhdeniye}



\begin{quote}
Before the new one: say the Russian for to burn oneself, then the Russian for to cut oneself.
\end{quote}

\subsection*{You'll want to know: \ru{рождение}}
\textbf{\ru{рождение}} (\emph{rozhdéniye}) — \textquotedblleft{}birth\textquotedblright{}.

\begin{quote}
\ru{Рождение} \ru{сына} \ru{изменило} \ru{их} \ru{жизнь}.
\end{quote}

\begin{quote}
\emph{Rozhdéniye sýna izmenílo ikh zhizn.}
\end{quote}

\begin{quote}
The birth of their son changed their life.
\end{quote}

Neuter. \emph{\ru{День} \ru{рождения}} — a birthday.

\begin{cousinweb}
\emph{\ru{рождение}} (EtymDB: \textquotedblleft{}birth\textquotedblright{}) is built on \emph{\ru{родить}} \textquotedblleft{}to give birth\textquotedblright{}, inherited from Proto-Slavic \emph{roditi} (EtymDB), with the Church Slavonic \emph{-\ru{жд}-} where native Russian has \emph{-\ru{ж}-}.
\end{cousinweb}

\begin{grammarlens}[title={What you've built}]
One more word for telling, explaining and describing at B1.
\end{grammarlens}

\subsection*{Your turn}
\begin{itemize}
\raggedright
  \item \emph{Say it:} \emph{rozhdéniye}
  \item \emph{Say it:} \emph{rozhdéniye}, once more
  \item \emph{Say it:} a sentence of your own with \emph{rozhdéniye}
  \item \emph{Recall:} say the Russian for to burn oneself, then the Russian for to cut oneself, then say \emph{rozhdéniye} again
\end{itemize}

\begin{tcolorbox}[breakable,colback=teal!4,colframe=teal!35!black,title={Before you move on}]
What does \ru{рождение} mean? (\textquotedblleft{}Birth\textquotedblright{}.) Say it once more.
\end{tcolorbox}
\section[\texorpdfstring{\ru{рождаться} (rozhdátsya)}{рождаться (rozhdátsya)}]{\ru{рождаться} (rozhdátsya) — to be born}

\label{lesson:RU-C303-rozhdatsya}



\begin{quote}
Before the new one: say the Russian for to cut oneself, then the Russian for birth.
\end{quote}

\subsection*{You'll want to know: \ru{рождаться}}
\textbf{\ru{рождаться}} (\emph{rozhdátsya}) — \textquotedblleft{}to be born\textquotedblright{}.

\begin{quote}
\ru{Каждой} \ru{весной} \ru{у} \ru{нашей} \ru{кошки} \ru{рождаются} \ru{котята}.
\end{quote}

\begin{quote}
\emph{Kázhdoy vesnóy u náshey kóshki rozhdáyutsya kotyáta.}
\end{quote}

\begin{quote}
Every spring our cat has kittens.
\end{quote}

Imperfective; perfective \emph{\ru{родиться}} (\emph{rodítsya}): \emph{\ru{я} \ru{родился} \ru{в} \ru{мае}} — I was born in May.

\begin{cousinweb}
\emph{\ru{рождаться}} is built on \emph{\ru{родить}} \textquotedblleft{}to give birth\textquotedblright{}, inherited from Proto-Slavic \emph{roditi} (EtymDB), in the Church Slavonic shape \emph{-\ru{жд}-}, with the reflexive \emph{-\ru{ся}}.
\end{cousinweb}

\begin{grammarlens}[title={What you've built}]
One more word for telling, explaining and describing at B1.
\end{grammarlens}

\subsection*{Your turn}
\begin{itemize}
\raggedright
  \item \emph{Say it:} \emph{rozhdátsya}
  \item \emph{Say it:} \emph{rozhdátsya}, once more
  \item \emph{Say it:} a sentence of your own with \emph{rozhdátsya}
  \item \emph{Recall:} say the Russian for to cut oneself, then the Russian for birth, then say \emph{rozhdátsya} again
\end{itemize}

\begin{tcolorbox}[breakable,colback=teal!4,colframe=teal!35!black,title={Before you move on}]
What does \ru{рождаться} mean? (\textquotedblleft{}To be born\textquotedblright{}.) Say it once more.
\end{tcolorbox}
\section[\texorpdfstring{\ru{младенец} (mladénets)}{младенец (mladénets)}]{\ru{младенец} (mladénets) — a baby, an infant}

\label{lesson:RU-C303-mladenets}



\begin{quote}
Before the new one: say the Russian for birth, then the Russian for to be born.
\end{quote}

\subsection*{You'll want to know: \ru{младенец}}
\textbf{\ru{младенец}} (\emph{mladénets}) — \textquotedblleft{}a baby, an infant\textquotedblright{}.

\begin{quote}
\ru{Младенец} \ru{спал} \ru{в} \ru{кроватке}.
\end{quote}

\begin{quote}
\emph{Mladénets spal v krovátke.}
\end{quote}

\begin{quote}
The baby was asleep in the cot.
\end{quote}

Genitive \emph{\ru{младенца}}. Bookish; in speech \emph{\ru{малыш}} or \emph{\ru{ребёнок}} (\emph{rebyónok}).

\begin{cousinweb}
\emph{\ru{младенец}} is a Church Slavonic word, \emph{mladĭnĭcĭ} (EtymDB), from Proto-Slavic \emph{moldŭ} \textquotedblleft{}young\textquotedblright{}; the native Russian shape of the same root is \emph{\ru{молодой}}.
\end{cousinweb}

\begin{grammarlens}[title={What you've built}]
One more word for telling, explaining and describing at B1.
\end{grammarlens}

\subsection*{Your turn}
\begin{itemize}
\raggedright
  \item \emph{Say it:} \emph{mladénets}
  \item \emph{Say it:} \emph{mladénets}, once more
  \item \emph{Say it:} a sentence of your own with \emph{mladénets}
  \item \emph{Recall:} say the Russian for birth, then the Russian for to be born, then say \emph{mladénets} again
\end{itemize}

\begin{tcolorbox}[breakable,colback=teal!4,colframe=teal!35!black,title={Before you move on}]
What does \ru{младенец} mean? (\textquotedblleft{}A baby, an infant\textquotedblright{}.) Say it once more.
\end{tcolorbox}
\section[\texorpdfstring{\ru{беременная} (berémennaya)}{беременная (berémennaya)}]{\ru{беременная} (berémennaya) — pregnant}

\label{lesson:RU-C303-beremennaya}



\begin{quote}
Before the new one: say the Russian for to be born, then the Russian for a baby.
\end{quote}

\subsection*{You'll want to know: \ru{беременная}}
\textbf{\ru{беременная}} (\emph{berémennaya}) — \textquotedblleft{}pregnant\textquotedblright{}.

\begin{quote}
\ru{Сестра} \ru{сейчас} \ru{беременна}.
\end{quote}

\begin{quote}
\emph{Sestrá seychás berémenna.}
\end{quote}

\begin{quote}
My sister is pregnant now.
\end{quote}

The feminine form, since it describes a woman; short form \emph{\ru{беременна}}.

\begin{cousinweb}
\emph{\ru{беременная}} (EtymDB: \textquotedblleft{}pregnant\textquotedblright{}) is usually traced to \emph{\ru{беремя}} \textquotedblleft{}a burden, an armful\textquotedblright{} (EtymDB: Old East Slavic \emph{beremę}), whose Church Slavonic shape is \emph{\ru{бремя}}: carrying a load. EtymDB lists no parent for \emph{\ru{беременная}}.
\end{cousinweb}

\begin{grammarlens}[title={What you've built}]
One more word for telling, explaining and describing at B1.
\end{grammarlens}

\subsection*{Your turn}
\begin{itemize}
\raggedright
  \item \emph{Say it:} \emph{berémennaya}
  \item \emph{Say it:} \emph{berémennaya}, once more
  \item \emph{Say it:} a sentence of your own with \emph{berémennaya}
  \item \emph{Recall:} say the Russian for to be born, then the Russian for a baby, then say \emph{berémennaya} again
\end{itemize}

\begin{tcolorbox}[breakable,colback=teal!4,colframe=teal!35!black,title={Before you move on}]
What does \ru{беременная} mean? (\textquotedblleft{}Pregnant\textquotedblright{}.) Say it once more.
\end{tcolorbox}
\section[\texorpdfstring{\ru{рожать} (rozhát)}{рожать (rozhát)}]{\ru{рожать} (rozhát) — to give birth}

\label{lesson:RU-C303-rozhat}



\begin{quote}
Before the new one: say the Russian for a baby, then the Russian for pregnant.
\end{quote}

\subsection*{You'll want to know: \ru{рожать}}
\textbf{\ru{рожать}} (\emph{rozhát}) — \textquotedblleft{}to give birth\textquotedblright{}.

\begin{quote}
\ru{Она} \ru{родила} \ru{дочку} \ru{в} \ru{марте}.
\end{quote}

\begin{quote}
\emph{Oná rodilá dóchku v márte.}
\end{quote}

\begin{quote}
She gave birth to a daughter in March.
\end{quote}

Imperfective; perfective \emph{\ru{родить}} (\emph{rodít}).

\begin{cousinweb}
\emph{\ru{рожать}} is inherited through Old East Slavic \emph{ražati} from Proto-Slavic \emph{radjati} (EtymDB), the imperfective partner of \emph{roditi} \textquotedblleft{}to give birth\textquotedblright{}.
\end{cousinweb}

\begin{grammarlens}[title={What you've built}]
One more word for telling, explaining and describing at B1.
\end{grammarlens}

\subsection*{Your turn}
\begin{itemize}
\raggedright
  \item \emph{Say it:} \emph{rozhát}
  \item \emph{Say it:} \emph{rozhát}, once more
  \item \emph{Say it:} a sentence of your own with \emph{rozhát}
  \item \emph{Recall:} say the Russian for a baby, then the Russian for pregnant, then say \emph{rozhát} again
\end{itemize}

\begin{tcolorbox}[breakable,colback=teal!4,colframe=teal!35!black,title={Before you move on}]
What does \ru{рожать} mean? (\textquotedblleft{}To give birth\textquotedblright{}.) Say it once more.
\end{tcolorbox}
